\section{Native operators}

This section records the compact operator layer required for the calculus. All operators below are discrete bookkeeping on declared traces, active cuts, frames/windows, and inherited support.

\subsection{Q4 chart and closure operators}\label{subsec:q4-closure-operators}
\begin{definition}[Q4 chart and closure operators]
\label{def:q4-closure-operator}
Given the axis-lift chart policy of \Cref{choice:axis-lift}, the Q4 chart operator records signed single-axis toggles on the canonical chart, the derived level index $\lambda=L_{\pm}/4$, the primitive beat-closure witness $L_{\pm}=2$, the beat-knot / cage-knot split, the minimal cage-band witness $(4{:}6)$, and the first nondegenerate enclosure as a four-cycle witness. The chart basis consists of vertices, Hamming-1 edges, square faces, cubes, and the 4-cell body; two-cycle returns are minimal returns and four-cycle enclosures are the first nondegenerate closure class on that basis.
\end{definition}

\paragraph{Depends on.} \Cref{choice:q4-chart,choice:axis-lift,def:q4-basis,def:cycle-typing,def:closure-predicate,def:beat-cage-knot,def:cage-band}.
\paragraph{Local substructure.} Axis-lift chart policy; Q4 cell basis; beat-knot / cage-knot typing; minimal cage-band $(4{:}6)$; minimal four-cycle enclosure.
\paragraph{Derivation sketch.} From the declared chart policy and the inherited support law, active ternary steps map to single-axis toggles on Q4. The closure predicate yields beat-knot versus cage-knot classes on the same chart, the minimal two-knot cage report is typed by $(4{:}6)$, and the first nondegenerate enclosure in this cubical basis is a four-cycle.

\subsection{Kernel-clock reconciliation, native loss channel, and rest channel}\label{subsec:kernel-clock-rest}
\begin{definition}[Kernel-clock reconciliation operator]
\label{def:kernel-clock-reconciliation}
Let $S_t$ be the declared local wavefront state, $A_t(S_t)$ the admissible successor set, and $P_t(\cdot\mid S_t)$ the declared local choice kernel on that set. Let $\beta_t$ be the backlog witness and let $F_t$ be the forced-dwell predicate. The realized update law is the reconciliation operator
\[
\widetilde P_t(\tau\mid S_t,\beta_t):=
\begin{cases}
\delta_0(\tau), & \text{if }F_t\text{ holds},\\[4pt]
\dfrac{P_t(\tau\mid S_t)}{\sum_{\sigma\in A_t(S_t)}P_t(\sigma\mid S_t)}, & \text{if }\tau\in A_t(S_t)\text{ and }F_t\text{ does not hold},\\[8pt]
0, & \text{otherwise.}
\end{cases}
\]
The kernel expresses local choice; backlog and forced dwell express clock-side debt service. The latter do not reparameterize the former.
\end{definition}

\begin{definition}[Derived biz rest rate]
\label{def:rest-biz}
With a declared attenuation accumulator $A_B(M;\omega)$ and burden reference $B_0:=1\,\mathrm{bip}$, define the derived internal rest rate by
\[
R_{\mathrm{rest}}^{(\mathrm{biz})}(M;\omega):=\frac{T_\Gamma(M;\omega)}{B_0}\,A_B(M;\omega)
\qquad [\mathrm{biz}].
\]
The primitive remains $R_{\mathrm{rest}}=u^{(\mathrm{loss})}$ in trit/bip.
\end{definition}

\paragraph{Depends on.} \Cref{choice:tau-readout,def:active-dwell,def:mt,def:loss-channel,def:backlog-reconciliation}.
\paragraph{Local substructure.} Admissibility count $m_t$; forced dwell; $\Gamma_{\mathrm{loss,raw}}$; $\Gamma_{\mathrm{loss}}$; $T_\Gamma$; primitive trit/bip channel; derived biz channel.
\paragraph{Derivation sketch.} From the declared gating rule and the active/dwell split, admissibility restriction yields $m_t<3$ and therefore positive loss on constrained ticks. Averaging on the cited active-cut family gives the native rest channel; the biz rate is the corresponding internal rate on the same ledger.

\subsection{Wavelet transport, trit carriage, and saturation}\label{subsec:wavelet-transport-saturation}
\begin{definition}[Wavelet transport operator]
\label{def:wavelet-transport-operator}
A wavelet transport operator packages a finite closure-carrying Memory-path segment together with its ordered cutwise field quantities, traverser/closure saturation, and chained trit carriage on the same active-cut family.
\end{definition}

\begin{definition}[Wavelet mode and transport witnesses]
\label{def:wavelet-mode-witness}
For a wavelet segment $e$ on a cited window $\omega$, let
\[
A(e;\omega):=N_+(e;\omega)+N_-(e;\omega),\quad Z(e;\omega):=N_0(e;\omega),\quad S(e;\omega):=N_+(e;\omega)-N_-(e;\omega),
\]
\[
E(e;\omega):=A(e;\omega)+Z(e;\omega)=|\omega\cap e|.
\]
The native ledgers are $(A,Z,S,E)$. Derived reporting witnesses are
\[
\kappa_{\mathrm{prog}}:=\frac{A}{E},\qquad c_{\mathrm{seat}}:=\frac{Z}{E}=1-\kappa_{\mathrm{prog}},\qquad \xi:=\frac{S}{E}.
\]
These are dictionary/rendering compressions only; the integer ledgers remain authoritative.
\end{definition}

\begin{definition}[Coupling exposure and propensity]
\label{def:coupling-exposure}
For a declared wavelet class $s$ on a cited readout family $\omega$, let $N_{\mathrm{exp}}(s;\omega)$ be the exposure count and $N_{\mathrm{cpl}}(s;\omega)$ the counted couplings attributed to those exposures. The empirical coupling propensity is
\[
p_{\mathrm{cpl}}(s;\omega):=\frac{N_{\mathrm{cpl}}(s;\omega)}{\max(1,N_{\mathrm{exp}}(s;\omega))}.
\]
This is the native counting contract used on the cited ledger.
\end{definition}

\paragraph{Depends on.} \Cref{def:ao-field,def:wavelet,def:saturation,def:traverser-closure-rule,def:loss-channel}.
\paragraph{Local substructure.} Wavelet segment; saturation; chained trit carriage; progression and coupling-seat witnesses; exposure/coupling counts.
\paragraph{Derivation sketch.} From the declared ternary readout and the wavelet/saturation ledgers on Memory paths. The visible stochastic layer is bookkeeping over the same cutwise transport operators.

\subsection{Push/pull, effective charge, and export/radiation operators}\label{subsec:pushpull-export-operators}
\begin{definition}[Push/pull boundary bookkeeping]
\label{def:push-pull-operator}
Push/pull is the chirality-resolved boundary-current bookkeeping on the cited readout family. With a fixed sign convention,
\[
J_\alpha(M;\omega):=\sum_{t\in\omega}\sum_{(i,j)\in\partial R(M;\omega)} \phi_{ij}(t)\,\mathbf 1\{\chi(t)=\alpha\},
\]
\[
J_\Omega(M;\omega):=\sum_{t\in\omega}\sum_{(i,j)\in\partial R(M;\omega)} \phi_{ij}(t)\,\mathbf 1\{\chi(t)=\Omega\}.
\]
The same readout supports the effective-charge proxy
\[
q_{\mathrm{eff}}(M;\omega):=\mathcal N\!\bigl(J_\alpha(M;\omega)-J_\Omega(M;\omega)\bigr)
\]
and the phase-lag witness $\Delta\Phi_{\Omega|\alpha}(M;\omega)$ under the declared pairing policy.
\end{definition}

\begin{definition}[Radiating/export window and emitted packet]
\label{def:radiation-operator}
A window $\omega$ is radiating/exporting for motif $M$ when net boundary export on $\partial R(M;\omega)$ is nonzero beyond tolerance and/or the boundary mutual-information witness exceeds the declared shell bound. The emitted/exported packet witness is the wavelet packet crossing the declared cut on that window.
\end{definition}

\paragraph{Depends on.} \Cref{def:push-pull-current,def:qeff,def:phase-lag,def:radiating-window,def:wavelet-mode-witness,def:coupling-exposure}.
\paragraph{Local substructure.} Push/pull sign convention; chirality-resolved boundary totals; effective-charge proxy; export witness; emitted packet.
\paragraph{Derivation sketch.} From the same oriented boundary-current histories used for lag, torsion, and shell/export classification. No external electrodynamic or thermodynamic law is invoked.

\subsection{Torsional curl and torsional residual}\label{subsec:torsional-operators}
\begin{definition}[Torsional curl and residual]
\label{def:torsional-operator}
Let $R(M;\omega)$ be the declared Region extractor for motif $M$ on the cited active-cut family. With a fixed boundary orientation, define the chirality-resolved boundary circulation
\[
\mathrm{Curl}_{\chi}(M;\omega):=
\sum_{t\in\omega}\sum_{(i,j)\in\partial R(M;\omega)}
\operatorname{sgn}_{ij}\,J^{\chi}_{ij}(t),
\qquad \chi\in\{\alpha,\Omega\},
\]
and the torsional curl difference
\[
\mathrm{Curl}^{g}(M;\omega):=\mathrm{Curl}_{\alpha}(M;\omega)-\mathrm{Curl}_{\Omega}(M;\omega).
\]
The torsional residual is
\[
T_{\mathrm{tors}}(M;\omega):=
\sum_{t\in\omega}\sum_{(i,j)\in\partial R(M;\omega)}
\operatorname{sgn}_{ij}\bigl(J^{\alpha}_{ij}(t)-J^{\Omega}_{ij}(t)\bigr)
= \mathrm{Curl}^{g}(M;\omega).
\]
By linearity of the AFC Stokes pairing, the same quantity has the region-side form
\[
\sum_{t\in\omega}\sum_{i\in R(M;\omega)}
\operatorname{div}\!\bigl(J^{\alpha}(t)-J^{\Omega}(t)\bigr)(i)
= T_{\mathrm{tors}}(M;\omega).
\]
In the symmetric zero-lag limit,
\[
J^{\alpha}=J^{\Omega}\Longrightarrow T_{\mathrm{tors}}(M;\omega)=0.
\]
\end{definition}

\paragraph{Depends on.} \Cref{def:chirality,def:torsional,def:phase-lag,def:backlog-reconciliation}.
\paragraph{Local substructure.} Chirality-resolved boundary currents; torsional curl; boundary-side residual; region-side equivalent; zero-lag limit.
\paragraph{Derivation sketch.} From the chirality-resolved boundary histories on the same active-cut family and the AFC Stokes pairing between boundary and region accounting.
\begin{remarkx}
This is the native Stokes-style witness for unresolved torsional imbalance on the cited readout family.
\end{remarkx}

\subsection{Cadence and polyrotor extractors}\label{subsec:cadence-polyrotor-operators}
\begin{definition}[Cadence extractor and shell-resolved accumulators]
\label{def:cadence-operator}
For a declared shell band or active-cut family, the cadence extractor returns cycle counts $N_k^{\mathrm{cyc}}(\omega)$ together with the shell-signature class and shell-resolved phase and rotor accumulators
\[
\Phi_k(\omega):=\sum_{t\in\omega}\Delta\Phi_k(t),
\qquad
R_k(\omega):=\sum_{t\in\omega}\Delta R_k(t),
\]
and drift witnesses
\[
\delta\Phi_{ij}(\omega):=\Phi_i(\omega)-\Phi_j(\omega),
\qquad
\delta R_{ij}(\omega):=R_i(\omega)-R_j(\omega).
\]
The corresponding cadence witness is
\[
\nu_k(\omega):=\frac{N_k^{\mathrm{cyc}}(\omega)}{|\omega|},
\]
so the reported cadence spectrum is $\mathcal N(M;\omega)=\{\nu_1(\omega),\ldots,\nu_K(\omega)\}$.
\end{definition}

\paragraph{Depends on.} \Cref{def:polyrotor,def:cadence-extractor,def:drift,def:cadence-bstar,def:shell-cadence-interface}.
\paragraph{Local substructure.} Cadence extractor; phase/rotor accumulators; drift witnesses; cadence spectrum; dominant residual-duration class.
\paragraph{Derivation sketch.} From the declared shell-band extraction and phase/rotor ledgers on the same cited readout family.
\begin{remarkx}
On shell-qualified windows, recurrent closure and phase-locking may be read as standing wavelet configurations on the cited support. Their unresolved phase disagreement is reported through backlog, torsional residual, cadence shift, and the native rest-channel on that same window. The cut-to-cut identity is unchanged; the effective window-level progress rate need not match branch-specific traversal duration, because part of the cited window is spent servicing $\alpha/\Omega$ mismatch through bridge and dwell burden. The dominant residual-duration class $B^*$ records the leading support class of that burden on the cited window under the declared extraction policy.
\end{remarkx}

\subsection{\texorpdfstring{$B^*$}{B*} ladder and histogram operators}\label{subsec:bstar-ladder-operators}
\begin{definition}[$B^*$ ladder and histogram operators]
\label{def:bstar-operator}
Given the return-duration extraction of \Cref{def:b-extraction}, the weighted histogram
\[
H_B(b;M,\omega):=\sum_{\kappa\in\mathrm{Occ}(M;\omega)} w(\kappa)\,\mathbf 1\{b(\kappa)=b\}
\]
induces the dominant residual-duration class
\[
B^*(M;\omega):=\arg\max_b H_B(b;M,\omega)
\]
and the support mass on that class
\[
\Lambda_b(M;\omega):=H_B(B^*(M;\omega);M,\omega).
\]
Under the declared minimal-pivot identity,
\[
b=2e+f,
\qquad f\in\mathbb Z_{\ge1},
\]
the ladder records the integer residual-duration burden on the cited window.
\end{definition}

\begin{definition}[Jensen discipline for the ladder]
\label{def:jensen-discipline}
Histogram and attenuation summaries must be formed on the declared sample set before downstream compression. If $T_{\mathrm{eval}}=n_1+n_2+n_3$ and
\[
\bar m:=\frac{n_1+2n_2+3n_3}{T_{\mathrm{eval}}},
\]
then the exact native loss-average is
\[
\frac{n_1+n_2\log_3\!\left(\frac{3}{2}\right)}{T_{\mathrm{eval}}},
\]
not the post-averaged surrogate $\log_3(3/\bar m)$. Equivalently, the Jensen gap
\[
J_{\mathrm{gap}}(\omega):=
\frac{n_1+n_2\log_3\!\left(\frac{3}{2}\right)}{T_{\mathrm{eval}}}
-
\log_3\!\left(\frac{3}{\bar m}\right)\ge 0
\]
must be respected on the cited admissibility ledger. Likewise, one must not replace histogram- or attenuation-level quantities such as
\[
\sum_b H_B(b;M;\omega)\,b^{-\widehat\Lambda(M;\omega)}
\]
by expressions built from averaged $b(\kappa)$ unless the aggregation policy explicitly licenses that reduction.
\end{definition}

\begin{definition}[Dominant-rung attenuation]
\label{def:attenuation-operator}
Using the same stack ledger $S(M;\omega)$ and residual-duration samples $b(\kappa)$ already required to compute $B^*(M;\omega)$, define the exact return attenuation accumulator
\[
A_B^{\mathrm{exact}}(M;\omega):=
\prod_{\kappa\in S(M;\omega)}\left(\frac{b(\kappa)}{B_0}\right)^{-L_{\pm}(\kappa)}\left(\frac{B^*(M;\omega)}{B_0}\right)^{-\delta_{\mathrm{core}}(M;\omega)}.
\]
Let the stack-load exponent be
\[
\widehat\Lambda(M;\omega):=\sum_{\kappa\in S(M;\omega)}L_{\pm}(\kappa)+\delta_{\mathrm{core}}(M;\omega).
\]
Under the declared dominant-rung approximation $b(\kappa)\approx B^*(M;\omega)$ on the cited stack,
\[
A_B(M;\omega)=\left(\frac{B^*(M;\omega)}{B_0}\right)^{-\widehat\Lambda(M;\omega)}.
\]
Thus $\Lambda_b$ remains the support mass on the dominant class, while $\widehat\Lambda$ is the stack-load exponent used in attenuation.
\end{definition}

\paragraph{Depends on.} \Cref{def:b-extraction,def:bstar-histogram,choice:hard-freeze,rem:jensen,def:cycle-typing}.
\paragraph{Local substructure.} Outbound traversal burden; return/closure burden; bridge/dwell service; residual-duration histogram; dominant class; support mass; shell-signature class; Jensen discipline; attenuation accumulator.
\paragraph{Derivation sketch.} From the declared traversal, closure, and bridge/dwell burdens one forms $b(\kappa)$, then $H_B$, $B^*$, and $\Lambda_b$ on the cited window; attenuation follows by algebra on the stack ledger.

\subsection{Finite path sums, least reconciliation, ternary interaction gate, and diagrammatic renderings}\label{subsec:pathsum-diagram-operators}
\begin{definition}[Finite path sum]
\label{def:path-sum}
A finite path sum is a finite sum over declared Memory paths, with each term given by a finite product of ledger-derived factors such as admissibility counts, chirality tags, phase factors, coupling factors, and attenuation weights.
\end{definition}

\begin{definition}[Least reconciliation]
\label{def:least-reconciliation}
Least reconciliation is the discrete minimization of a declared cost functional over admissible trace families. The cost may include loss burden, dwell cost, boundary residual, and return-debt service on the same active-cut family.
\end{definition}

\begin{definition}[Ternary interaction gate]
\label{def:ternary-gate}
The primitive wavelet tick alphabet is $\tau\in\{-1,0,+1\}$. The balanced ternary gate is defined by identifying $\{-1,0,+1\}$ with $\mathbb Z_3$ via $\chi(-1)=2$, $\chi(0)=0$, $\chi(+1)=1$ and setting
\[
a\oplus_3 b:=\chi^{-1}\!\bigl((\chi(a)+\chi(b))\bmod 3\bigr).
\]
This implements cancellation/cascade bookkeeping such as $+1\oplus_3(-1)=0$ and $+1\oplus_3(+1)=-1$.
\end{definition}

\begin{definition}[Collision couplings and diagrammatic renderings]
\label{def:collision-diagram}
Push/pull collision couplings of duon/duad and tetron/tetrad classes are represented as native bookkeeping on declared traces. Cut diagrams, closure diagrams, path-family diagrams, and finite diagrammatic renderings package the same ledgers in compact report form.
\end{definition}

\paragraph{Depends on.} \Cref{def:wavelet,def:field-families,def:observer-coupling,def:push-pull-operator,def:attenuation-operator}.
\paragraph{Local substructure.} Finite path sums; least reconciliation; ternary interaction gate; collision couplings; diagrammatic renderings; export packets.
\paragraph{Derivation sketch.} From finite sums and products over declared traces on the same AO field and wavelet ledgers. Diagrammatic language packages those same ledgers in a compact report form.
\begin{remarkx}
Finite path sums and diagrammatic renderings record the same cut bookkeeping on the cited window; they are report forms of the native line. They package already-declared AO field, wavelet, boundary-current, torsional, cadence, and return-ladder ledgers.
\end{remarkx}


\subsection{Field persistence and lifecycle}\label{subsec:field-persistence-lifecycle}
\begin{definition}[Field persistence and lifecycle]
\label{def:field-lifecycle-operator}
For a declared motif id $M$ and cited window family $\omega$, the persistence operator records the birth tick, death tick, TTL hazard, censoring status, rupture type, and any product/by-product census reported after rupture. In native units,
\[
\mathrm{TTL}_{\mathrm{bip}}(M;\omega)=t_{\mathrm{death}}-t_{\mathrm{birth}}.
\]
When TTL is used in histograms or ladder weighting, the weighting rule must be declared on the same manifest.
\end{definition}

\paragraph{Depends on.} \Cref{def:ttl-hazard,def:decay-census,def:field-families}.
\paragraph{Local substructure.} TTL hazard; rupture predicate; right-censoring; decay; by-product census.
\paragraph{Derivation sketch.} Pure integer/tick bookkeeping on the same motif-conditioned ledgers under the declared persistence contract.

\subsection{Freezing, rebalance, transmutation, and dynamic recruitment}\label{subsec:freeze-rebalance-recruitment}
\begin{definition}[Freezing, confinement witness, and rebalance operators]
\label{def:freezing-operator}
Freezing is the reported sustained admissibility restriction on a cited motif-conditioned stream. The confinement witness is the operator-side realization of the declared confinement predicate on the same shell-qualified window and shell-signature class. Rebalance dwell is the explicit debt-service dwell attached to emission, absorption, transmutation, or recruitment events. Dynamic volume recruitment is the growth of the effective closure range when freezing persists.
\end{definition}

\paragraph{Depends on.} \Cref{def:freeze-recruit,def:mt,def:backlog-reconciliation,def:drift,def:shell-signature,def:confinement-predicate}.
\paragraph{Local substructure.} Cross-shell drift; backlog witness; forced dwell; shell-signature class; $B^*$-anchored persistence; confinement witness; failure criterion; transmutation / recruitment of outer closure range.
\paragraph{Derivation sketch.} On the cited readout family,
\[
\text{cross-shell drift}\Rightarrow \beta_t>0\Rightarrow \text{forced dwell}\Rightarrow m_t<3\Rightarrow R_{\mathrm{rest}}>0.
\]
When the same shell-qualified window also satisfies the declared proximity and $B^*$-dominance condition, the window is reported as confined. Failure of that witness triggers the declared transmutation event operator and, if sustained, dynamic recruitment of the outer closure range.

\subsection{Rest-channel derivation}\label{subsec:native-rest-derivation}
\begin{definition}[Rest-channel derivation]
\label{def:native-rest-derivation}
The rest-facing chain is the declared operator composition
\[
m_t\rightarrow u_t^{(\mathrm{loss})}\rightarrow \Gamma_{\mathrm{loss}}\rightarrow B^*\rightarrow \widehat\Lambda\rightarrow A_B\rightarrow R_{\mathrm{rest}}\rightarrow R_{\mathrm{rest}}^{(\mathrm{biz})}.
\]
Here
\[
u_t^{(\mathrm{loss})}:=\log_3\!\left(\frac{3}{m_t}\right),
\qquad
U^{(\mathrm{loss})}(M;\omega):=\sum_{t\in(\omega\cap M)_{\mathrm{act}}} u_t^{(\mathrm{loss})},
\]
\[
R_{\mathrm{rest}}(M;\omega)=u^{(\mathrm{loss})}(M;\omega):=\frac{1}{T_{\mathrm{eval}}(M;\omega)}U^{(\mathrm{loss})}(M;\omega)
=\log_3\!\bigl(\Gamma_{\mathrm{loss}}(M;\omega)\bigr)\quad[\mathrm{trit/bip}],
\]
with
\[
\Gamma_{\mathrm{loss}}(M;\omega):=\left(\prod_{t\in(\omega\cap M)_{\mathrm{act}}}\frac{3}{m_t}\right)^{1/T_{\mathrm{eval}}(M;\omega)},
\qquad
T_\Gamma(M;\omega):=\Gamma_{\mathrm{loss}}(M;\omega)-1,
\]
\[
b(\kappa)=e_{\alpha}(\kappa)+e_{\Omega}(\kappa)+f(\kappa),
\qquad
B^*(M;\omega)=\arg\max_b H_B(b;M;\omega),
\]
\[
\widehat\Lambda(M;\omega):=\sum_{\kappa\in S(M;\omega)}L_{\pm}(\kappa)+\delta_{\mathrm{core}}(M;\omega),
\qquad
A_B(M;\omega)=\left(\frac{B^*(M;\omega)}{B_0}\right)^{-\widehat\Lambda(M;\omega)},
\]
and the derived internal rate is
\[
R_{\mathrm{rest}}^{(\mathrm{biz})}(M;\omega):=\frac{T_\Gamma(M;\omega)}{B_0}\,A_B(M;\omega)
\quad[\mathrm{biz}].
\]
\end{definition}

\paragraph{Depends on.} \Cref{def:mt,def:loss-channel,def:b-extraction,def:bstar-histogram,def:attenuation-operator,choice:hard-freeze,rem:jensen}.
\paragraph{Local substructure.} Admissibility count; native loss channel; return-duration ledger; dominant rung; stack load; attenuation accumulator; primitive trit/bip channel; derived biz rate.
\paragraph{Derivation sketch.} Direct algebra on the cited admissibility and return ledgers. The zero-lag/unconstrained limit $m_t=3$ yields $R_{\mathrm{rest}}=0$ and $T_\Gamma=0$.

\subsection{Momentum-like native proxy and rest-frame windows}\label{subsec:momentum-proxy}
\begin{definition}[Momentum-like native proxy]
\label{def:momentum-proxy-operator}
Using the normalized speed $v(\omega)$ on a cited window, define the momentum-like native proxy by
\[
p_{\mathrm{prox}}(M;\omega):=E_{\mathrm{tot}}(M;\omega)\,v(\omega),
\]
where $E_{\mathrm{tot}}$ is the declared total biz-valued rate assembled from the native rest-facing channel and any declared additional native totals on the same ledger. This is the native proxy channel assembled on the cited ledger.
\end{definition}

\paragraph{Depends on.} \Cref{def:rest-frame-window,def:native-rest-derivation}.
\paragraph{Local substructure.} Normalized speed; rest-frame window class; native proxy channel.
\paragraph{Derivation sketch.} From axis-traversal speed on the cited window and the already-declared native rate channels.

\subsection{Radiation and export dynamics}\label{subsec:radiation-export-dynamics}
\begin{definition}[Radiation and export dynamics]
\label{def:radiation-dynamics-operator}
For a motif-conditioned Region extractor $R(M;\omega)$ with boundary $\partial R(M;\omega)$, define shell qualification by
\[
\mathrm{shell}(M;\omega)\iff \mathrm{net\_export}(M;\omega)=0\text{ within tolerance }\wedge MI_{\partial}(M;\omega)\le \varepsilon_I,
\]
and define the radiating/export class by
\[
\mathrm{rad}(M;\omega)\iff \mathrm{net\_export}(M;\omega)\neq 0\ \vee\ MI_{\partial}(M;\omega)>\varepsilon_I.
\]
The emission/export packet witness is the wavelet packet whose boundary crossing certifies the export on the same cited ledgers.
\end{definition}

\paragraph{Depends on.} \Cref{def:radiating-window,def:push-pull-operator,def:wavelet,def:coupling-exposure}.
\paragraph{Local substructure.} Motif$\to$Region extractor; boundary MI export; shell-qualified window; radiating/export window; packet witness; shell falsifier.
\paragraph{Derivation sketch.} From the same native boundary ledgers used for shell and export classification on the cited window.

\subsection{Family dynamics and subclass rules}\label{subsec:family-dynamics-subclass}
\begin{definition}[Family dynamics and subclass rules]
\label{def:family-dynamics-operator}
The public field families carry declared dynamic constraints in addition to their naming keys: allowed occupancy states, closure-lock conditions, saturation behavior by subclass, and admissible transition rules. For duad instances, the lifetime observable is often time-to-first-interaction rather than intrinsic dissolution time; this caveat must be preserved whenever duad persistence is reported.
\end{definition}

\paragraph{Depends on.} \Cref{def:field-families,def:traverser-closure-rule,def:saturation,def:field-lifecycle-operator}.
\paragraph{Local substructure.} Family persistence behavior; structural key; seat/traversal constraints; closure-lock conditions; subclass saturation behavior; admissible transition rules.
\paragraph{Derivation sketch.} Port of the native naming/lifecycle constraints into a compact operator-facing contract.

\subsection{Collision and export-window calculus}\label{subsec:collision-export-calculus}
\begin{definition}[Collision and export-window calculus]
\label{def:collision-export-operator}
A collision window is a declared export window on which incoming wavelet ledgers, the ternary interaction gate, and any production/export tick are all recorded on the same cited trace family. On a cited short export window, the same enclosure and shell-signature classes used for confinement and transmutation also type the collision bookkeeping, so the packet witness and residual check are read on the same family-conditioned ledgers. The native residual check is
\[
\mathcal R_{\mathrm{cut}}:=S_{\mathrm{out}}-\bigl(S_{\mathrm{in}}+S_{\mathrm{prod}}\bigr),
\]
which must vanish within tolerance on the same cut ledger. Exchange rules are implemented with the balanced ternary gate $\oplus_3$ on the same declared collision/export ledger.
\end{definition}

\paragraph{Depends on.} \Cref{def:ternary-gate,def:collision-diagram,def:radiation-dynamics-operator}.
\paragraph{Local substructure.} Collision window as export window; XOR exchange; production/export tick; packet witness; residual check.
\paragraph{Derivation sketch.} From native tick ledgers and boundary bookkeeping on the declared collision/export window. Class-pair interaction tables, when recorded, follow the same native ledger and the same declared collision/export-window contract.
